\section{Linear Programming: The Simplex Method}

Questions that came up during the book club session:
\begin{enumerate}
    \item Does convexity of a problem imply that the KKT conditions are sufficient?
          Yes, it does.
          We begin with a useful definition.
          \begin{Def}[Convex problem]
              Given the functions $c_i, i \in \mathcal{I}$, $c_j, j \in \mathcal{E}$ and $f$, in the problem \eqref{eq:constrained-problem}, we say that the optimization problem is \textit{convex} if\footnote{
                  This definition would be less awkward had the feasibility condition been the more conventional $c_i(x) \leq 0$. It would then define $c_i$ to be convex instead of concave.
              }
              \begin{align*}
                   & \text{All }c_i\text{ are concave.}\quad \implies \quad \nabla c_i(x)^T(y - x) \geq c_i(y) - c_i(x) \\
                   & \text{All }c_j\text{ are affine.}\quad \implies \quad \nabla c_j(x)^T(y - x) = c_j(y) - c_j(x)     \\
                   & f\text{ is convex.}\quad \implies \quad \nabla f(x)^T(y - x) \leq f(y) - f(x)
              \end{align*}
              \label{def:convex-problem}
          \end{Def}

          \begin{Thm}[Sufficiency of KKT for convex problems]
              Given the problem \eqref{eq:constrained-problem}, and a point $x^*\in\Omega$ for which the KKT conditions are satisfied for some $\lambda^*$.
              If the problem is convex as defined in definition \ref{def:convex-problem}, then $x^*$ is a global minimizer.
          \end{Thm}
          \textit{Proof} Pick any $x\in\Omega$.
          From the definition \ref{def:convex-problem}, we have for the active constraints that
          \begin{align*}
              \nabla c_i(x^*)^T(x - x^*) & \geq c_i(x) - c_i(x^*) = c_i(x) \geq 0 \\
              \nabla c_j(x^*)^T(x - x^*) & = c_j(x) - c_j(x^*) = 0                \\
          \end{align*}
          Then, by inspecting the convexity condition on the objective function we obtain
          \begin{align*}
              f(x) - f(x^*) \geq \nabla f(x^*)^T(x - x^*) = \sum_{i\in\mathcal{A}(x^*)} \lambda_i \nabla c_i (x^*)^T (x - x^*)  \geq 0
          \end{align*}
          Thus, $f(x) \geq f(x^*)$ for any $x \in \Omega$. \hfill$\blacksquare$
    \item Can we formulate some more compelling, detailed version of the proof of part (ii) of the strong duality theorem?
          It turns out we can, and it is based on the basic logical relation, that $A \implies B \iff \lnot B \implies \lnot A$.
          An example in words could be ''\textit{A number divisible by four is even implies that no odd number is divisible by four}''.
          Recalling the weak duality theorem,
          \begin{Thm}[Weak Duality]
              If $x$ is a primal feasible point and $\lambda$ is a dual feasible point, then
              \[
                  b^T\lambda \leq c^T x
              \]
          \end{Thm}
          we now state the second part of the strong duality theorem,
          \begin{Thm}[Strong Duality II \cite{mangasarian1994nonlinear}]
              If the primal problem is unbounded, the dual problem is infeasible.
              Conversely, if the dual problem is unbounded, the primal is infeasible.
          \end{Thm}
          \textit{Proof}
          The statement
          $$
              \text{The primal is unbounded } \implies \text{ The dual is infeasible}
          $$
          is equivalent to
          $$
              \text{The dual is feasible } \implies \text{ The primal is bounded}.
          $$
          So, to show that unboundedness of the dual implies that the primal is infeasible, $\Omega = \emptyset$, it suffices to show that feasibility of the dual implies that the primal is bounded, which follows from the weak duality theorem.
          Every dual feasible point $\lambda$ provides a lower bound for the primal optimal objective via
          $$
              b^T\lambda \leq c^T x, \quad \forall \> x\in X
          $$
          A symmetric argument can be made for the case
          $$
              \text{The dual is unbounded } \implies \text{ The primal is infeasible}
          $$
          \hfill$\blacksquare$

\end{enumerate}

\subsection{Solutions}
The exercises chosen for this chapter are 13.1,
\subsubsection{13.1}
To convert this problem to standard form we need to do the following
\begin{itemize}
    \item Reformulate it into a minimization problem. Done by negating the objective function.
    \item Convert inequality constraints to equality constraints.
    \item Replace unbounded variables by non-negative variables.
\end{itemize}
The first point is self-explanatory.
\[
    \operatorname{max} c^T x + d^T y \to \operatorname{min} -c^T x - d^T y
\]
For the second point we begin by splitting the bound on $y$ into two separate constraints
\[
    l \leq y \leq u \to \begin{array}{l}
        y \leq u \\
        y \geq l
    \end{array}
\]
So the full set of constraints are
\begin{align*}
    A_1 x         & = b_1    \\
    A_2 x + B_2 y & \leq b_2 \\
    y             & \leq u   \\
    y             & \geq l
\end{align*}
Adding slack (and \textit{surplus} for the lower bounds) variables yields the following picture
\begin{align*}
    A_1 x               & = b_1  \\
    A_2 x + B_2 y + s_1 & = b_2  \\
    y + s_2             & = u    \\
    y - s_3             & = l    \\
    s_1, s_2, s_3       & \geq 0
\end{align*}
Finally, we split the currently unbounded variables $x$ and $y$ into negative and positive parts as $x \to x^+ - x^-$ with $x^+, x^-\in \mathbb{R}^+$ and similarly for $y$.
This gives the final problem
\[
    \operatorname{min} -c^T\left(x^+ - x^-\right) - d^T\left(y^+ - y^-\right)
\]
subject to
\begin{align*}
    A_1 \left(x^+ - x^-\right)                                    & = b_1   \\
    A_2 \left(x^+ - x^-\right) + B_2 \left(y^+ - y^-\right) + s_1 & = b_2   \\
    \left(y^+ - y^-\right) + s_2                                  & = u     \\
    \left(y^+ - y^-\right) - s_3                                  & = l     \\
    x^+, x^-, y^+, y^-, s_1, s_2, s_3                             & \geq 0.
\end{align*}
One can formulate the equality constraints as a single matrix equation by constructing a large matrix out of the different blocks above:
\[
    \begin{pmatrix}
        A_1 & -A_1 & 0   & 0    & 0 & 0 & 0  \\
        A_2 & -A_2 & B_2 & -B_2 & I & 0 & 0  \\
        0   & 0    & I   & -I   & 0 & I & 0  \\
        0   & 0    & I   & -I   & 0 & 0 & -I \\
    \end{pmatrix}\begin{pmatrix}
        x^+ \\
        x^- \\
        y^+ \\
        y^- \\
        s_1 \\
        s_2 \\
        s_3
    \end{pmatrix} =
    \begin{pmatrix}
        b_1 \\
        b_2 \\
        u   \\
        l
    \end{pmatrix}
\]
\subsubsection{13.2}
Verify that the dual of
\[
    \operatorname{max} b^T\lambda, \quad \text{subject to } A^T\lambda \leq c
\]
is the primal.
Recall that the dual problem is:
\[
    \underset{\lambda \geq 0}{\operatorname{max}} \> \underset{x}{\operatorname{min}} \> \mathcal{L}(x, \lambda)
\]
The Lagrangian is, in this case (note the order of the variables and the negation of $b^T\lambda$)
\begin{align*}
    \mathcal{L}(\lambda, x) & = -b^T\lambda - x^T(c - A^T\lambda) \\
                            & = (x^T A^T - b^T)\lambda - x^T c.
\end{align*}
The dual of this problem is then
\[
    \underset{x \geq 0}{\operatorname{max}} \> \underset{\lambda}{\operatorname{min}} \> \mathcal{L}(\lambda, x)
\]
Unless $b^T - x^T A^T = 0 \implies Ax = b$ is satisfied, the problem is unbounded. Therefore we may enforce $Ax = b$ as a constraint, and we obtain the problem
\[
    \operatorname{max} -x^Tc, \text{ Subject to } Ax = b, \quad x \geq 0 ,
\]
which can be turned into the primal by changing $\operatorname{max} -x^Tc \to \operatorname{min} c^Tx$.
